\documentclass[10pt,a4paper]{amsart}
\usepackage[margin=19mm]{geometry}
\usepackage{amsmath,amssymb,mathtools}
\usepackage[scr=rsfs,cal=euler]{mathalfa}
\usepackage{xfrac}
\usepackage[unicode,hidelinks,bookmarks=false]{hyperref}
\newcommand{\ie}{i.e.\ }
\newcommand{\cf}{cf.\ }
\newcommand{\resp}{respectively\ }
\newcommand{\Subsection}{\subsection}
\providecommand{\nobreakdash}{\nobreak\hbox{-}}
\setlength{\emergencystretch}{2em}
\multlinegap0pt
\usepackage{bm}
\usepackage{bbm}
\usepackage{dsfont}
\usepackage{xspace}
\usepackage{cancel}
\usepackage{mathtools}
\usepackage{amsthm}
\makeatletter
\@ifundefined{theorem}{\newtheorem{theorem}{Theorem}}{}
\makeatother
\newcommand{\nn}{\mathbb{N}}
\newcommand{\zz}{\mathbb{Z}}
\title[The extensible no-$(k(n)+1)$-in-line problem]{The extensible no-$(k(n)+1)$-in-line problem}
\author{Tam\'as G\'abriel, M\'at\'e J\'anosik, D\'avid Melj\'an, Benedek N\'ador}
\date{Source: arXiv:2606.02843v1 (CC BY 4.0)}
\begin{document}
\maketitle

\noindent\emph{Source and licence.} The extensible no-$(k(n)+1)$-in-line problem. Version arXiv:2606.02843v1 (CC BY 4.0), distributed under the Creative Commons Attribution 4.0 International licence, as recorded in the arXiv licence metadata for this exact version and shown on the abstract page https://arxiv.org/abs/2606.02843v1. Full rights notice: licenses/arxiv-2606.02843-CC-BY-4.0.txt.\par\smallskip

\noindent\textit{Editorial scope.} Counted unit: the Theorem whose source label is \texttt{None} in the article (source0.tex lines 947--951, source label \texttt{None}), delivered together with the article's own complete original proof of it (source0.tex lines 953--1001, one \texttt{proof} environment of 49 lines). No mathematics is written, repaired or re-derived: statement and proof are the article's own text line for line. Required context retained uncounted: none. Every definition, display and label of those lines is retained unchanged. Omitted from this package and not redistributed: 1--946, 952--952, 1002--1291. Nothing inside a retained excerpt is omitted or altered: every retained body is delivered exactly as the article's lines are written, so no omission receipt of any kind accompanies this package. Citations: the retained excerpts carry no citation command at all, so no bibliography entry of the article is transcribed and no reference is redistributed. Reference adaptation: none was needed and none was made; every reference inside the delivered unit resolves inside it, so no unresolved reference or citation is produced. Printed slips kept literal: no printed word, symbol, hypothesis or number of the counted statement or of its proof was altered. Layout and class: the article's own class is \texttt{article} with packages including amssymb, amsthm, thmtools, amsmath, tcolorbox, pgfplots, babel, inputenc, hyperref, cleveref, tabularx, array, enumitem, wrapfig; the wrapper is the lane's pinned \texttt{amsart} editorial wrapper at 10pt on A4, which restates the theorem-like environments the retained text needs and adds the article's own macro definitions that the retained text needs (\texttt{\textbackslash nn}, \texttt{\textbackslash zz}), with extra wrapper packages (bm, bbm, dsfont, xspace, cancel, mathtools). The delivered numbering is the unit's own. Assets: no retained span contains a figure environment, a graphics-inclusion command, a PDF-page inclusion, a table or a TikZ picture; no image file is copied into this package, the package declares no asset dependency and no image file is redistributed with it. Licence: version arXiv:2606.02843v1 of this article is distributed under the Creative Commons Attribution 4.0 International licence, as recorded in the arXiv licence metadata for this exact version and shown on the abstract page https://arxiv.org/abs/2606.02843v1. A full rights notice is delivered at licenses/arxiv-2606.02843-CC-BY-4.0.txt. Characters: every retained excerpt is pure ASCII, so the delivered engine, XeLaTeX with its default Latin Modern text font, reports no missing character.
\begin{theorem}
Let $k:\nn^+\to\mathbb{R}^+$ be a monotone increasing function and $S\subseteq\zz^{2}$ a set of lattice points such that $\left|S\cap\ell\cap[1,n]^{2}\right|\le k(n)$ for all lines $\ell$. If $S_n:=S\cap [1,n]^2$ satisfies
$$\liminf_{n\to\infty} \frac{|S_n|}{k(n)\cdot n} \ge 0.897,$$
then $k(n) = \Omega( n^c)$, where $c$ is a positive absolute constant.
\end{theorem}

\begin{proof}
By the condition on $S_n$, there is a constant $n_0$ satisfying $\frac{|S \cap [1,n]^2|}{n k(n)} >1-\varepsilon$ for every $n \ge n_0$, where $\varepsilon=0.1036$. Pick $n_0$ such that it further satisfies $k(n_0) > 0$. Subdivide $[1, mn_0]^2$ into $m\times m$ tiles of size $n_0\times n_0$. 
 
 By the definition of $n_0$, there are at least $(1-\varepsilon)(m-1)n_0k((m-1)n_0)$ points from $S$ in $[1,(m-1)n_0]^2$. $[1,(m-1)n_0]\times[1,mn_0]$ contains $(m-1)n_0$ disjoint vertical lines inside the  square $[1, mn_0]^2$. Consequently, it contains at most $(m-1)n_0k(mn_0)$ points of $S$.
  Therefore, the rectangle $[1,(m-1)n_0]\times[(m-1)n_0+1,mn_0]$ contains at most $$(m-1)n_0k(mn_0)-(1-\varepsilon)(m-1)n_0k((m-1)n_0)$$ points of $S$.  
  
  The rectangle $[1,mn_0]\times[(m-1)n_0+1,mn_0]$ is the disjoint union of $n_0$ horizontal lines inside the square $[1, mn_0]^2$. The other $(m-1)n_0$ horizontal lines contain at most $(m-1)n_0k(mn_0)$ points inside the square $[1, mn_0]^2$. By the definition of $n_0$, the square $[1, mn_0]^2$ contains at least $(1-\varepsilon)mn_0k(mn_0)$ points from $S$. Therefore the rectangle $[1,mn_0]\times[(m-1)n_0+1,mn_0]$ contains at least $(1-m\varepsilon)n_0k(mn_0)$ points from $S$. Together with the previous bound, this implies that $[(m-1)n+1,mn]^2$ contains at least 

  $$   (1-m\varepsilon)n_0k(mn_0)-\big((m-1)n_0k(mn_0)-(1-\varepsilon)(m-1)n_0k((m-1)n_0)\big)=$$ 
 $$=((m-1-(m-1)\varepsilon)k((m-1)n_0)-(m-2+m\varepsilon)k(mn_0))n_0$$
 points of $S$. This bound applies to all squares in the main diagonal of $[1, mn_0]^2$.

The $m$ squares of size $n_0\times n_0$ in the main diagonal of $[1,mn_0]^2$ can be covered by $2n_0-1$ diagonal lines of slope 1, inside the square $[1,mn_0]^2$, which means they contain at most $(2n_0-1)k(mn_0)\le2n_0k(mn_0)$ points from $S$. Comparing this bound with the previous lower bound on these squares yields 
$$\sum_{i=1}^{m}((i-1-(i-1)\varepsilon)k((i-1)n_0)-(i-2+i\varepsilon)k(in_0))n_0\le2n_0k(mn_0)$$
$$\sum_{i=1}^{m-1}(2-2i\varepsilon)k(in_0)-(m-2+m\varepsilon)k(mn_0)\le 2k(mn_0)$$
\begin{equation}
\sum_{i=1}^{m-1}2(1-i\varepsilon)k(in_0)\le m(1+\varepsilon)k(mn_0).
\label{eq:sum_bound}
\end{equation}

In the following we recursively establish an inequality between $k(n_0)$ and $k(5n_0)$.

First apply inequality \ref{eq:sum_bound} for $m=2$:

$$2(1-\varepsilon)k(n_0)\le2(1+\varepsilon)k(2n_0)$$
$$\frac{1-\varepsilon}{1+\varepsilon}k(n_0)\le k(2n_0).$$

Proceeding with $m=3$ yields

$$2(1-\varepsilon)k(n_0)+2(1-2\varepsilon)k(2n_0)\le3(1+\varepsilon)k(3n_0).$$
Using $\frac{1-\varepsilon}{1+\varepsilon}k(n_0)\le k(2n_0)$:
$$\left((1-\varepsilon)+(1-2\varepsilon)\frac{1-\varepsilon}{1+\varepsilon}\right)2k(n_0)\le 3(1+\varepsilon)k(3n_0)$$
$$\frac{2(1-\varepsilon)(2-\varepsilon)}{1+\varepsilon}k(n_0)\le 3(1+\varepsilon)k(3n_0)$$
$$\frac{2(1-\varepsilon)(2-\varepsilon)}{3(1+\varepsilon)^2}k(n_0)\le k(3n_0).$$
Applying inequality \ref{eq:sum_bound} for $m=4$ yields
$$2(1-\varepsilon)k(n_0)+2(1-2\varepsilon)k(2n_0)+2(1-3\varepsilon)k(3n_0)\le4(1+\varepsilon)k(4n_0)$$
$$\left(\frac{2(1-\varepsilon)(2-\varepsilon)}{1+\varepsilon}+2(1-3\varepsilon)\frac{2(1-\varepsilon)(2-\varepsilon)}{3(1+\varepsilon)^2}\right)k(n_0)\le4(1+\varepsilon)k(4n_0)$$
$$\frac{2(1-\varepsilon)(2-\varepsilon)(5-3\varepsilon)}{3(1+\varepsilon)^2}k(n_0)\le 4(1+\varepsilon)k(4n_0)$$
$$\frac{(1-\varepsilon)(2-\varepsilon)(5-3\varepsilon)}{6(1+\varepsilon)^3}k(n_0)\le k(4n_0).$$
For $m=5$ the process yields
$$2(1-\varepsilon)k(n_0)+2(1-2\varepsilon)k(2n_0)+2(1-3\varepsilon)k(3n_0)+2(1-4\varepsilon)k(4n_0)\le5(1+\varepsilon)k(5n_0)$$
$$\left(\frac{2(1-\varepsilon)(2-\varepsilon)(5-3\varepsilon)}{3(1+\varepsilon)^2}+2(1-4\varepsilon)\frac{(1-\varepsilon)(2-\varepsilon)(5-3\varepsilon)}{6(1+\varepsilon)^3}\right)k(n_0)\le 5(1+\varepsilon)k(5n_0)$$
$$\frac{(1-\varepsilon)(2-\varepsilon)(5-3\varepsilon)(3-2\varepsilon)}{15(1+\varepsilon)^4}\le k(5n_0).$$
$\varepsilon=0.1036$ means $Ck(n_0)\le k(6n_0)$, where $C=1.00054>1$ is a constant. We can repeat this argument for every $n\ge n_0$, including every $n=5^an_0$, where $a\in\nn$, this result gives us 
$$C^ak(n_0)\le k(5^an_0)=k(n)$$
$$C^{\log_5\frac{n}{n_0}}k(n_0)=5^{\log_5C\log_5\frac{n}{n_0}}k(n_0)=\left(\frac{n}{n_0}\right)^{\log_5C}k(n_0)\le k(n).$$
Setting $c=\log_5C\approx 0.0003>0$ and the monotonicity of $k(n)$ means $k(n)=\Omega(n^c)$.
    
\end{proof}

\end{document}
